% questions.tex —— 2008 年·普通高等学校招生考试·天津卷·文科 · 题目内容片段（无答案）
% 由 tools/generate_projects.py 依来源总账自 2008/2008_数学真题_天津文_试题.md 生成，勿手改（改源 MD 后重跑）


\gkpart{一、选择题}

\begin{question}
\gkti{1}设集合 $U=\{x\in\mathbf N\mid0<x\leqslant8\}$，$S=\{1,2,4,5\}$，$T=\{3,5,7\}$，则 $S\cap(\complement_UT)=$（　）

（A）$\{1,2,4\}$

（B）$\{1,2,3,4,5,7\}$

（C）$\{1,2\}$

（D）$\{1,2,4,5,6,8\}$

\end{question}
\begin{question}
\gkti{2}设变量 $x$，$y$ 满足约束条件 $\begin{cases}x-y\geqslant0,\\ x+y\leqslant1,\\ x+2y\geqslant1,\end{cases}$ 则目标函数 $z=5x+y$ 的最大值为（　）

（A）$2$

（B）$3$

（C）$4$

（D）$5$

\end{question}
\begin{question}
\gkti{3}函数 $y=1+\sqrt x$（$0\leqslant x\leqslant4$）的反函数是（　）

（A）$y=(x-1)^2$（$1\leqslant x\leqslant3$）

（B）$y=(x-1)^2$（$0\leqslant x\leqslant4$）

（C）$y=x^2-1$（$1\leqslant x\leqslant3$）

（D）$y=x^2-1$（$0\leqslant x\leqslant4$）

\end{question}
\begin{question}
\gkti{4}若等差数列 $\{a_n\}$ 的前 $5$ 项和 $S_5=25$，且 $a_2=3$，则 $a_7=$（　）

（A）$12$

（B）$13$

（C）$14$

（D）$15$

\end{question}
\begin{question}
\gkti{5}设 $a$，$b$ 是两条直线，$\alpha$，$\beta$ 是两个平面，则 $a\perp b$ 的一个充分条件是（　）

（A）$a\perp\alpha$，$b\parallel\beta$，$\alpha\perp\beta$

（B）$a\perp\alpha$，$b\perp\beta$，$\alpha\parallel\beta$

（C）$a\subset\alpha$，$b\perp\beta$，$\alpha\parallel\beta$

（D）$a\subset\alpha$，$b\parallel\beta$，$\alpha\perp\beta$

\end{question}
\begin{question}
\gkti{6}把函数 $y=\sin x$（$x\in\mathbf R$）的图象上所有的点向左平行移动 $\dfrac{\pi}3$ 个单位长度，再把所得图象上所有点的横坐标缩短到原来的 $\dfrac12$ 倍（纵坐标不变），得到的图象所表示的函数是（　）

（A）$y=\sin\left(2x-\dfrac{\pi}3\right)$，$x\in\mathbf R$

（B）$y=\sin\left(\dfrac x2+\dfrac{\pi}6\right)$，$x\in\mathbf R$

（C）$y=\sin\left(2x+\dfrac{\pi}3\right)$，$x\in\mathbf R$

（D）$y=\sin\left(2x+\dfrac{2\pi}3\right)$，$x\in\mathbf R$

\end{question}
\begin{question}
\gkti{7}设椭圆 $\dfrac{x^2}{m^2}+\dfrac{y^2}{n^2}=1$（$m>0$，$n>0$）的右焦点与抛物线 $y^2=8x$ 的焦点相同，离心率为 $\dfrac12$，则此椭圆的方程为（　）

（A）$\dfrac{x^2}{12}+\dfrac{y^2}{16}=1$

（B）$\dfrac{x^2}{16}+\dfrac{y^2}{12}=1$

（C）$\dfrac{x^2}{48}+\dfrac{y^2}{64}=1$

（D）$\dfrac{x^2}{64}+\dfrac{y^2}{48}=1$

\end{question}
\begin{question}
\gkti{8}已知函数 $f(x)=\begin{cases}x+2,&x\leqslant0,\\ -x+2,&x>0,\end{cases}$ 则不等式 $f(x)\geqslant x^2$ 的解集为（　）

（A）$[-1,1]$

（B）$[-2,2]$

（C）$[-2,1]$

（D）$[-1,2]$

\end{question}
\begin{question}
\gkti{9}设 $a=\sin\dfrac{5\pi}7$，$b=\cos\dfrac{2\pi}7$，$c=\tan\dfrac{2\pi}7$，则（　）

（A）$a<b<c$

（B）$a<c<b$

（C）$b<c<a$

（D）$b<a<c$

\end{question}
\begin{question}
\gkti{10}设 $a>1$，若对于任意的 $x\in[a,2a]$，都有 $y\in[a,a^2]$ 满足方程 $\log_ax+\log_ay=3$，这时 $a$ 的取值的集合为（　）

（A）$\{a\mid1<a\leqslant2\}$

（B）$\{a\mid a\geqslant2\}$

（C）$\{a\mid2\leqslant a\leqslant3\}$

（D）$\{2,3\}$

\end{question}
\gkpart{二、填空题}

\begin{question}
\gkti{11}一个单位共有职工 $200$ 人，其中不超过 $45$ 岁的有 $120$ 人，超过 $45$ 岁的有 $80$ 人。为了调查职工的健康状况，用分层抽样的方法从全体职工中抽取一个容量为 $25$ 的样本，应抽取超过 $45$ 岁的职工\gkblank{}人。

\end{question}
\begin{question}
\gkti{12}$\left(x+\dfrac2x\right)^5$ 的二项展开式中 $x^3$ 的系数为\gkblank{}。（用数字作答）

\end{question}
\begin{question}
\gkti{13}若一个球的体积为 $4\sqrt3\pi$，则它的表面积为\gkblank{}。

\end{question}
\begin{question}
\gkti{14}已知平面向量 $\boldsymbol a=(2,4)$，$\boldsymbol b=(-1,2)$，若 $\boldsymbol c=\boldsymbol a-(\boldsymbol a\cdot\boldsymbol b)\boldsymbol b$，则 $|\boldsymbol c|=$\gkblank{}。

\end{question}
\begin{question}
\gkti{15}已知圆 $C$ 的圆心与点 $P(-2,1)$ 关于直线 $y=x+1$ 对称。直线 $3x+4y-11=0$ 与圆 $C$ 相交于 $A$，$B$ 两点，且 $|AB|=6$，则圆 $C$ 的方程为\gkblank{}。

\end{question}
\begin{question}
\gkti{16}有 $4$ 张分别标有数字 $1$，$2$，$3$，$4$ 的红色卡片和 $4$ 张分别标有数字 $1$，$2$，$3$，$4$ 的蓝色卡片，从这 $8$ 张卡片中取出 $4$ 张卡片排成一行。如果取出的 $4$ 张卡片所标的数字之和等于 $10$，则不同的排法共有\gkblank{}种。（用数字作答）

\end{question}
\gkpart{三、解答题}

\begin{question}
\gkti{17}已知函数 $f(x)=2\cos^2\omega x+2\sin\omega x\cos\omega x+1$（$x\in\mathbf R$，$\omega>0$）的最小正周期是 $\dfrac{\pi}2$。

（1）求 $\omega$ 的值；

（2）求函数 $f(x)$ 的最大值，并且求使 $f(x)$ 取得最大值的 $x$ 的集合。

\end{question}
\begingroup
\renewenvironment{question}{\Needspace{0.46\textheight}}{}
\begin{question}
\gkti{18}甲、乙两个篮球运动员互不影响地在同一位置投球，命中率分别为 $\dfrac12$ 与 $p$，且乙投球 $2$ 次均未命中的概率为 $\dfrac1{16}$。

（1）求乙投球的命中率 $p$；

（2）求甲投球 $2$ 次，至少命中 $1$ 次的概率；

（3）若甲、乙两人各投球 $2$ 次，求两人共命中 $2$ 次的概率。

\Needspace{0.4\textheight}
\end{question}
\endgroup
\begin{question}
\gkti{19}如图，在四棱锥 $P-ABCD$ 中，底面 $ABCD$ 是矩形。已知 $AB=3$，$AD=2$，$PA=2$，$PD=2\sqrt2$，$\angle PAB=60^\circ$。

（1）证明 $AD\perp$ 平面 $PAB$；

（2）求异面直线 $PC$ 与 $AD$ 所成的角的大小；

（3）求二面角 $P-BD-A$ 的大小。
\nopagebreak[4]
\begin{center}% 2008 年天津卷文科第 19 题：四棱锥 P-ABCD，底面 ABCD 是矩形，
%   AB=3，AD=2，PA=2，PD=2√2，∠PAB=60°。
% 来源 PDF：gaokaomath-src/普通高考/2008/2008天津文.pdf
% SHA-256：4c8f1c998f95d805901f2e0d93fdbb5901ce4bc67fd7c3ca712083325837bb3d
% 原页：2008天津文.pdf 第 1 页（A3 三栏）第 3 栏，书页码 489。
% 重绘依据：原图底面为平行四边形 ABCD，A 为后侧顶点；P—B、P—C、P—D、
%   B—C、C—D 实线，P—A、A—B、A—D、B—D 虚线，A 处有 PA⊥AD 的直角标记。
%   坐标按原图 300dpi 像素位置等比换算（B(540,510)、C(825,515)、D(975,370)、
%   A(685,370)、P(640,170)）。
\begin{tikzpicture}[scale=1.0, semithick,
  every node/.style={inner sep=1.5pt}]
  \coordinate (B) at (0.00,0.00);
  \coordinate (C) at (2.85,0.00);
  \coordinate (D) at (4.35,1.45);
  \coordinate (A) at (1.45,1.45);
  \coordinate (P) at (1.00,3.45);

  \draw (P) -- (B);
  \draw (P) -- (C);
  \draw (P) -- (D);
  \draw (B) -- (C);
  \draw (C) -- (D);
  \draw[dashed] (P) -- (A);
  \draw[dashed] (A) -- (B);
  \draw[dashed] (A) -- (D);
  \draw[dashed] (B) -- (D);

  % A 处 PA⊥AD 的直角标记
  \coordinate (A1) at ($(A)!0.26cm!(P)$);
  \coordinate (A2) at ($(A)!0.26cm!(D)$);
  \draw (A1) -- ($(A1)+(A2)-(A)$) -- (A2);

  \node[below left=1pt] at (A) {$A$};
  \node[below left] at (B) {$B$};
  \node[below right] at (C) {$C$};
  \node[right] at (D) {$D$};
  \node[above] at (P) {$P$};
\end{tikzpicture}
\end{center}
\nopagebreak[4]
\end{question}
\begin{question}
\gkti{20}已知数列 $\{a_n\}$ 中，$a_1=1$，$a_2=2$，且 $a_{n+1}=(1+q)a_n-qa_{n-1}$（$n\geqslant2$，$q\neq0$）。

（1）设 $b_n=a_{n+1}-a_n$（$n\in\mathbf N^*$），证明 $\{b_n\}$ 是等比数列；

（2）求数列 $\{a_n\}$ 的通项公式；

（3）若 $a_3$ 是 $a_6$ 与 $a_9$ 的等差中项，求 $q$ 的值，并证明：对任意的 $n\in\mathbf N^*$，$a_n$ 是 $a_{n+3}$ 与 $a_{n+6}$ 的等差中项。

\end{question}
\begin{question}
\gkti{21}设函数 $f(x)=x^4+ax^3+2x^2+b$（$x\in\mathbf R$），其中 $a$，$b\in\mathbf R$。

（1）当 $a=-\dfrac{10}3$ 时，讨论函数 $f(x)$ 的单调性；

（2）若函数 $f(x)$ 仅在 $x=0$ 处有极值，求 $a$ 的取值范围；

（3）若对于任意的 $a\in[-2,2]$，不等式 $f(x)\leqslant1$ 在 $[-1,1]$ 上恒成立，求 $b$ 的取值范围。

\end{question}
\begin{question}
\gkti{22}已知中心在原点的双曲线 $C$ 的一个焦点是 $F_1(-3,0)$，一条渐近线的方程是 $\sqrt5x-2y=0$。

（1）求双曲线 $C$ 的方程；

（2）若以 $k$（$k\neq0$）为斜率的直线 $l$ 与双曲线 $C$ 相交于两个不同的点 $M$，$N$，线段 $MN$ 的垂直平分线与两坐标轴围成的三角形的面积为 $\dfrac{81}2$，求 $k$ 的取值范围。

\end{question}